\documentclass[10pt]{article}
\usepackage[a5paper, hmargin=14mm, top=8mm, bottom=12mm]{geometry}
\usepackage{amsmath, amssymb, amsthm}
\newtheorem{theorem}{Theorem}
\pagenumbering{gobble}
\title{Braiding Fibonacci anyons}
\author{Michael Atiyah \and Henri Poincar\'e \and Richard Feynman}
\date{}

\begin{document}
\maketitle

\section{Introduction}
Topological quantum computers keep qubits in
the fusion space of anyons~\cite{kitaev03}.
Fibonacci anyons have one nontrivial particle,
$\tau$, with the fusion rule
\[ \tau \otimes \tau = 1 \oplus \tau. \]
%TYPED%

\begin{theorem}
  The fusion space of $n$ copies of $\tau$
  has dimension $F_{n+1}$, a Fibonacci number.
\end{theorem}

\section{Braiding}
Exchanging two anyons that fuse to $c$ multiplies
the state by the phase $R^{\tau\tau}_c$, where
\[ R^{\tau\tau}_1 = e^{-4\pi i/5}, \qquad
   R^{\tau\tau}_\tau = e^{3\pi i/5}. \]
Changing the order of fusion is the $F$-move, and
the two satisfy the pentagon and hexagon
equations~\cite{kitaev06, wang10}. For how this
becomes a computer, see the survey~\cite{nayak08}.

\bibliographystyle{plain}
{\small\bibliography{references}}
\end{document}
